\paragraph{Baseline Conditions:}

Unless otherwise specified during the parameter sweeps, all numerical simulations evaluate the fuel-air mixtures at a baseline unburnt inlet temperature of $300$ \si{K} and an absolute pressure of $1$ \si{atm}. The oxidizer is consistently modelled as standard air, approximated as a molar ratio of one part oxygen to 3.76 parts nitrogen ($O_2 + 3.76 N_2$).

\subsection{Complete combustion}

This section evaluates the theoretical limits of methane and hydrogen combustion by assuming 100\% complete, one way reactions. The variations in heat capacity $C_p$ and adiabatic flame temperature $T_{ad}$ are analysed without the influence of molecular dissociation.

\paragraph{a)}

\autoref{fig:cp vs t} shows the variation in the mass specific heat $C_p$ for $\ce{H_2}$, $\ce{O_2}$, $\ce{N_2}$, $\ce{CO_2}$ and $\ce{CH_4}$ with temperature for a $T$ range of 300 to 2400 \si{K}. The figure uses a logarithmic scale for the $C_p$ axis due to the large discrepancies between the different substances.

\begin{figure}[H]
    \centering
    \includegraphics[width=0.5\linewidth]{Figures/1a_variation_cp_temp_cantera_logarithmic.pdf}
    \caption{Variation of the Mass Specific Heat Capacity with Temperature}
    \label{fig:cp vs t}
\end{figure}



\autoref{fig:cp_variations_comparison} shows the variation in the mixture specific heat capacity $C_{p,\text{mix}}$ for the reactants and combustion products assuming the reactants to be air and either $\ce{H_2}$ or $\ce{CH_4}$, for a range of equivalence ratios between 0.5 and 1.4. 


\begin{figure}[H]
    \centering
    % Left Subfigure: Methane
    \begin{subfigure}[b]{0.48\textwidth}
        \centering
        \includegraphics[width=\textwidth]{Figures/CH4_mixure_new_legend.png}
        \caption{Methane ($\ce{CH_4}$)-Air mixture.}
        \label{fig:cp_methane}
    \end{subfigure}
    \hfill % Adds horizontal space between the two figures
    % Right Subfigure: Hydrogen
    \begin{subfigure}[b]{0.48\textwidth}
        \centering
        \includegraphics[width=\textwidth]{Figures/H2_mixure_new_legend.png}
        \caption{Hydrogen ($\ce{H_2}$)-Air mixture.}
        \label{fig:cp_hydrogen}
    \end{subfigure}
    
    % Main Caption and Label for cross-referencing
    \caption{Variation of frozen mixture specific heat capacity ($C_{p,\text{mix}}$) for reactants and combustion products across a temperature range of 300\,K to 2400\,K at 1\,atm. } 
    \label{fig:cp_variations_comparison}
\end{figure}
%\todo{mudar loc da legenda}



\paragraph{b)} In all cases from \autoref{fig:cp_variations_comparison}, $C_{p,\text{mix}}$ changes with temperature consistently with the theory. As $T$ rises, higher molecular vibrational energy modes become excited, allowing the gas molecules to absorb more thermal energy per degree of temperature increase. 

As $\phi$ increases, the $C_p$ lines for the reactants shift upwards. This is expected for richer mixtures, as both $\ce{H_2}$ and $\ce{CH_4}$ have much higher mass-specific heat capacities than air, in its majority composed by $\ce{N_2}$ and $\ce{O_2}$, replacing air with fuel in the mixture makes the overall $C_p$ increase. When comparing both fuels, the $\ce{H_2}$ and air mixtures display significantly higher $C_p$ values compared to the $\ce{CH_4}$ air mixtures. This can be attributed to $\ce{H_2}$'s low molecular mass, resulting in an increase in the mass specific heat capacity. 

When comparing the reactants and products, for $\ce{H_2}$ air mixtures, the reactants have a higher value of $C_p$ than the products, since the combustion consumes the $\ce{H_2}$ gas, which has a high heat capacity. On the other hand, for $\ce{CH_4}$ and air mixtures, the $C_p$ of the reactants grows more steeply with temperature, often crossing above the products, since methane combustion produces polyatomic species such as $\ce{H_2O}$ and $\ce{CO_2}$ which possess several different vibrational modes, causing their heat capacity to rise rapidly at higher temperatures. 

%\todo{reduce word count to 200}

%part C
\paragraph{c)} \autoref{fig:1c fig numerical} plots the flame adiabatic temperature for the complete combustion of $\ce{H_2}$ and air for the range of equivalence ratios [0.5, 0.8, 1.0, 1.2, 1.4]. 

\begin{figure}[H]
    \centering
    \includegraphics[width=0.5\linewidth]{Figures/1c_adiabatic_flame_temp_eq_ratio_fixed.pdf}
    \caption{Adiabatic Flame Temperature for complete combustion of $\ce{H_2}$ and air, as well as $\ce{CH_4}$ and air}
    \label{fig:1c fig numerical}
\end{figure}

%\todo{review complete combustion}

% o antigo The graph highlights several relevant trends. Firstly, it can be seen that the adiabatic flame temperature does not peak at a stoichiometric ratio, with $\phi=1.0$, instead, this quantity peaks on the richer equivalence ratios. For methane, this peak happens at $\phi=1.1$, while hydrogen peaks much further into the rich region at $\phi=1.3$. The shift towards the richer side of the graph for the peaks happens since the excess of fuel suppresses highly endothermic dissociation reactions (such as $\ce{H2O -> OH + H}$). Since hydrogen has a low mass, excess $\ce{H2}$ acts as a very weak thermal diluent, allowing its peak to shift much further rich than methane. 

\autoref{fig:1c fig numerical} highlights several relevant trends for the thermodynamics of ideal complete combustion. 

Firstly, the adiabatic flame temperature $T_{ad}$ peaks at stoichiometry ($\phi=1.0$) for both fuels. Since the theoretical model assumes no endothermic dissociation, only stoichiometric combustion releases the maximum possible chemical energy. Any excess air or excess fuel act as a thermal diluent, absorbing heat and lowering the final temperature. 

Secondly, it can be observed that hydrogen burns significantly hotter than methane across the entire range. Complete methane combustion produces carbon dioxide. Since \ce{CO2} is a heavy polyatomic molecule, with multiple active vibrational modes, it possesses a substantially higher heat capacity than the \ce{H2O} and \ce{N2} produced by hydrogen. This causes \ce{CO2} to act as a stronger thermal sink, keeping the methane flame relatively cooler. 

Finally, the absolute temperatures are high, approaching 3200 \si{K} for \ce{H2}. This demonstrates how preheating reactants to 1100 \si{K} adds enthalpy to the flow, which is conserved and substantially elevates the final thermal state.

%\todo{A ver se percebi, o hydroge n aumenta tanto o cp?}
%Moreover, $\ce{H2}$ burns significantly hotter than $\ce{CH4}$ cross the entire range. This is due to the formation of $\ce{CO2}$ as a product of methane combustion, since carbon dioxide is a heavy polyatomic molecule with a large heat capacity, which absorbs thermal energy and keeps the methane flame relatively cooler. Finally, the data shows that preheating the reactants to 1100 K results in large peak adiabatic flame temperatures which approach 3000 K. This demonstrates how the pre-combustion enthalpy can drastically elevate the final thermal state.

%\todo{also make 200 words}

In conclusion, under ideal complete combustion, the absence of endothermic dissociation allows the flame temperature to peak exactly at stoichiometry, when $\phi = 1.0$. Furthermore, the high heat capacity of the \ce{CO2} produced during methane combustion acts as a thermal sink, rendering theoretical methane flames inherently cooler than hydrogen flames.

\subsection{Chemical equilibrium}

Moving beyond the ideal assumptions, this section uses Cantera's equilibrium solver to account for real high-temperature molecular dissociation. The true adiabatic flame temperatures are evaluated against variations in equivalence ratio, operating pressure, and inlet preheating.

\autoref{fig:numerical part 2a}, \autoref{fig:numerical part 2b} and \autoref{fig:numerical part 2c} respectively represent the variation in adiabatic flame temperature with equivalence ratio, with pressure, and with air inlet temperature for combustion of $\ce{H2}$ and $\ce{CH4}$ with air. Finally, \autoref{fig:numerical part 2d} represents the $\ce{NO}$ emissions in ppm for a range of different equivalence ratios.  

\begin{figure}[H]
    \centering
    \includegraphics[width=0.5\linewidth]{Figures/part2a.pdf}
    \caption{Adiabatic Flame Temperature for \ce{CH4} and \ce{H2} for different equivalence ratios at baseline conditions.}
    \label{fig:numerical part 2a}
\end{figure}

\autoref{fig:numerical part 2a}  introduces a real world chemical equilibrium, which accounts for the high temperature molecular dissociation. Unlike the ideal complete combustion model in \autoref{fig:1c fig numerical}, which peaked exactly at stoichiometry, the equilibrium temperatures peak on the slightly rich side, near $\phi=1.05$ for methane and $\phi=1.1$ for hydrogen. Since dissociation is a highly endothermic process that substantially lowers flame temperature, pushing the mixture slightly rich provides excess fuel that chemically suppresses these dissociation reactions, such as \ce{H2O -> OH + H}. By suppressing dissociation, the flame retains more thermal energy, causing the maximum adiabatic temperature to shift past $\phi = 1.0$.

%\autoref{fig:numerical part 2a} shows a shift in the adiabatic flame temperature peak towards lower equivalence ratios, compared to \autoref{fig:1c fig numerical}. In the current figure, the maximum $T_{ad}$ happens at $\phi = 1.0$ for $\ce{CH4}$ and at $\phi = 1.1$ for $\ce{H2}$. This earlier $\phi$ peak can be attributed to the lower inlet temperature of 300 K \  
%\todo{Erly as in, erly in phi? e isso é pq da inlet temp?}. 
%Since dissociation is an endothermic process that grows exponentially with temperature, the lower peak temperatures of around 2400 K make this phenomena much less pronounced, and as such they make the peak in adiabatic flame temperature closer to stoichiometric. 
%\todo{Rever isto pq da mudança do grafico anterior}

\begin{figure}[H]
    \centering
    \includegraphics[width=0.5\linewidth]{Figures/part2b.pdf}
    \caption{Variation of Adiabatic Flame Temperature for \ce{CH4} and \ce{H2} for Pressures at the baseline conditions.}
    \label{fig:numerical part 2b}
\end{figure}

\autoref{fig:numerical part 2b} highlights that an increase in pressure strictly increases the flame temperature, however, it is also apparent that this effect exhibits diminishing returns, as the increase in pressure from 10 to 40 bar yields only an increase in temperature of 8.64 K for $\ce{CH4}$ and 20.67 K for $H_2$. This behaviour can be explained by the lean condition at which the calculations are performed, since in this lower temperature regime of between 2000 and 2160 K, molecular dissociation is less prevalent compared to stoichiometric or rich conditions. Since dissociation is less pronounced, the effect of Le Chatelier's Principle, where the system shifts its chemical equilibrium towards the side with the fewer total moles in a reaction to counteract the pressure rise, is much more subtle, and as a result the temperature rise with an increase in pressure plateaus. 

\begin{figure}[H]
    \centering
    \includegraphics[width=0.5\linewidth]{Figures/part2c.pdf}
    \caption{Variation of Adiabatic Flame Temperature for \ce{CH4} and \ce{H2} for different Inlet Temperatures at the baseline conditions.}
    \label{fig:numerical part 2c}
\end{figure}
\autoref{fig:numerical part 2c} shows that preheating the incoming reactants increases the final flame temperature. However, this temperature rise is strictly less than the increase in inlet temperature ($\Delta T_{ad} < \Delta T_{in}$). This dampened response is driven by two factors: first, the specific heat capacity ($C_p$) of the gas mixture is significantly higher at elevated temperatures, meaning more thermal energy is required per degree of temperature rise. Second, in accordance with Le Chatelier's Principle, the higher peak temperatures drive the system to oppose the temperature rise by shifting its chemical equilibrium toward highly endothermic dissociation reactions. These reactions absorb a substantial portion of the added sensible enthalpy, further suppressing the final temperature.

\begin{figure}[H]
    \centering
    \includegraphics[width=0.5\linewidth]{Figures/part2d.pdf}
    \caption{\ce{NO} emissions in ppm at the baseline conditions for different equivalence ratios.}
    \label{fig:numerical part 2d}
\end{figure}

Finally, \autoref{fig:numerical part 2d} shows that $\ce{NO}$ emissions do not peak at a stoichiometric ratio of $\phi=1.0$, nor at the absolute temperature maximum ($\phi=1.1$ for $\ce{H2}$). Instead, this quantity is largest on the lean side, at $\phi=0.9$ for Methane and $\phi = 0.8$ for Hydrogen. For rich mixtures, $\ce{NO}$ emissions drop steeply for both fuels. This formation of $\ce{NO}$ is governed by the Extended Zeldovich Mechanism, which states that significant $\ce{NO}$ production requires two simultaneous conditions, which are extreme temperatures to break the highly stable $\ce{N2}$ triple bond, as well as an abundance of free oxygen to react with the nitrogen. This can create a thermodynamic balance across equivalence ratios. Although the flame reaches its maximum temperature at slightly rich conditions, the lack of free oxygen stops $\ce{NO}$ production. Consequently, $\ce{NO}$ emissions peak on the lean side, between $\phi = 0.8-0.9$, which is the optimal range where the flame remains at high temperature but still possesses enough unburnt oxygen to drive the Zeldovich reactions. 


Concluding, chemical equilibrium demonstrates that endothermic dissociation visibly limits peak flame temperatures, forcing the maximum $T_{ad}$ to shift to slightly fuel-rich mixtures. While increasing pressure and preheating reactants do elevate temperatures, their effectiveness is dampened at extreme conditions by Le Chatelier's principle and rising specific heat capacities.

%part a

\subsection{Estimating NOx}

This section investigates the formation of \ce{NOx} emissions, driven primarily by the temperature sensitive Extended Zeldovich Mechanism, across three distinct levels of modelling fidelity: algebraic global balances, infinite time chemical equilibrium, and finite rate 1D flame simulations.

To estimate the \ce{NO} mass fraction as a function of the equivalence ratio $\phi$, the exact number of residual \ce{N2} and \ce{O2} moles remaining after ideal complete combustion must first be determined. Rather than re-deriving the global mass balances, the previously established lean-to-stoichiometric global reactions (\autoref{eq:CH4_lean} for methane and \autoref{eq:H2_lean} for hydrogen) are utilized. 

From these established equations, the unreacted oxygen ($\nu_{\ce{O2},i}$) and inert nitrogen ($\nu_{\ce{N2},i}$) available for \ce{NO} formation per mole of fuel can be extracted directly:
\begin{itemize}
    \item \textbf{For Methane:} $\nu_{\ce{O2},i} = 2 \left(\frac{1}{\phi} - 1\right)$ and $\nu_{\ce{N2},i} = \frac{7.52}{\phi}$
    \item \textbf{For Hydrogen:} $\nu_{\ce{O2},i} = 0.5 \left(\frac{1}{\phi} - 1\right)$ and $\nu_{\ce{N2},i} = \frac{1.88}{\phi}$
\end{itemize}

After the amount of leftover \ce{O2} and \ce{N2} is found, this is subjected to the equilibrium adiabatic flame temperature $T_{eq}$ which was calculated in Part 2. At this temperature, the equilibrium reaction for the formation of \ce{NO} is described by \autoref{eq: equilibrium reaction 3a}, and the equilibrium constant $K_p$ can be calculated using the standard Gibbs free energy of reaction $\Delta G_{R} ^{\circ}$ at $T_{eq}$ (obtained through Cantera), as per \autoref{eq: gibbs free energy 3a}. 

\noindent
\begin{minipage}{0.5\textwidth}
\begin{equation}\label{eq: equilibrium reaction 3a}
    \ce{N2 + O2} \rightleftharpoons \ce{2 NO}
\end{equation}
\end{minipage}%
\begin{minipage}{0.5\textwidth}
\begin{equation}\label{eq: gibbs free energy 3a}
    K_p = \exp{\left( - \frac{\Delta G_{R} ^{\circ}}{R T_{eq}}\right)}
\end{equation}
\end{minipage}

Since the total number of moles does not change during this reaction (as $\Delta \nu = 2 -1 -1 = 0$), the pressure cancels out of the equilibrium equation. Defining $x$ as the change in the stoichiometric coefficient (the moles of \ce{O2} and \ce{N2} that react to form \ce{NO}), we can express the final equilibrium composition ($\nu_{eq}$) as per \autoref{eq: composition at equilibrium}. Substituting these coefficients into \autoref{eq: kp defenition} yields the evaluated definition of $K_p$. 
    
\noindent
\begin{minipage}{0.42\textwidth}
\begin{equation}\label{eq: composition at equilibrium}
\begin{aligned}
    \nu_{\ce{O2}, eq} &= \nu_{\ce{O2}, i} - x \\
    \nu_{\ce{N2}, eq} &= \nu_{\ce{N2}, i} - x \\
    \nu_{\ce{NO}, eq} &= 2x
\end{aligned}
\end{equation}
\end{minipage}%
\hfill
\begin{minipage}{0.55\textwidth}
\begin{equation}\label{eq: kp defenition}
    K_p = \frac{(\nu_{\ce{NO}, eq})^2}{(\nu_{\ce{N2}, eq})(\nu_{\ce{O2}, eq})} = \frac{(2x)^2}{(\nu_{\ce{N2},i}-x)(\nu_{\ce{O2},i}-x)}
\end{equation}
\end{minipage}

\begin{equation}\label{eq: 3a quadratic polynomial}
    (4-K_p)x^2 + K_p(\nu_{\ce{N2}, i} + \nu_{\ce{O2}, i})x - K_p (\nu_{\ce{N2}, i})(\nu_{\ce{O2}, i}) = 0
\end{equation}

Finally, expanding the denominator and moving all terms to one side yields the quadratic \autoref{eq: 3a quadratic polynomial}. The value for $x$ is taken to be the positive root of this polynomial. Once $x$ is solved, the moles of \ce{NO} produced per mole of fuel are found via $\nu_{\ce{NO}, eq} = 2x$. To plot the data, this molar amount is converted into a mass fraction $Y_{\ce{NO}}$ by multiplying by the molecular weight of \ce{NO} (30.01 \si{kg/kmol}) and dividing by the total initial mass of the reactant mixture for either the \ce{CH4} or \ce{H2} global reactions. \autoref{fig:part 3a} shows the graph resulting from this calculation for both fuels across all evaluated equivalence ratios. 

\begin{figure}[H]
    \centering
    \includegraphics[width=0.5\linewidth]{Figures/part3a.pdf}
    \caption{NO Mass Fraction as a function of $\phi$}
    \label{fig:part 3a}
\end{figure}
 


%part b


To compare the numerically computed results for both the 1D flames and equilibrium, \autoref{fig:comparison emissions} \footnote{Note: The hand-calculated NO emission for $\phi$ = 1.0 in plot (b) is exactly 0 ppm, which drops to the bottom axis boundary due to the logarithmic scale} shows the results for the \ce{NO} computation for question 3a, the equilibrium with CANTERA from 2d and the 1D flame simulation using CHEM1D. 

\begin{figure}[H]
    \centering
    
    % Left Figure
    \begin{subfigure}[b]{0.32\textwidth}
        \centering
        \includegraphics[width=\textwidth]{Figures/part2d_comparison_scale.pdf}
        \caption{\ce{NO} emissions from part 2d}
        \label{fig:2d emissions comparison}
    \end{subfigure}
    \hfill % Adds horizontal space
    % Center Figure
    \begin{subfigure}[b]{0.32\textwidth}
        \centering
        \includegraphics[width=\textwidth]{Figures/part3a no emissions ppm comparison scale.pdf}
        \caption{\ce{NO} emissions from part 3a}
        \label{fig:3a emissions comparison}
    \end{subfigure}
    \hfill % Adds horizontal space
    % Right Figure
    \begin{subfigure}[b]{0.32\textwidth}
        \centering
        \includegraphics[width=\textwidth]{Figures/Part3b_1D_Flame_NO_ppm_CHEM1D.pdf} % Replace filename
        \caption{\ce{NO} emissions from CHEM1D} % Replace caption
        \label{fig:chem 1D emissions comparison}
    \end{subfigure}
    
    \caption{Comparison of emissions from the \ce{NO} computation and equilibrium with CANTERA. (With adjusted graph bounds for comparison)}
    \label{fig:comparison emissions}
\end{figure}

The following comments can be done regarding both figures:

\begin{enumerate}[label=(\roman*)]
    \item At lean conditions, the two step computation and CANTERA equilibrium predict nearly identical \ce{NO} levels. This is because the massive bulk of unreacted \ce{O2} from lean combustion completely dominates the Zeldovich reaction, making any additional oxygen radicals from high temperature dissociation mathematically negligible.
    \item At stoichiometry ($\phi = 1.0$), the hand calculation predicts zero \ce{NO} because it assumes 100\% oxygen consumption. CANTERA correctly predicts high \ce{NO} since extreme temperatures force \ce{CO2} and \ce{H2O} to dissociate, supplying the trace free oxygen required to drive the Zeldovich mechanism.
    \item The Equilibrium models assume infinite residence time, allowing the slow Zeldovich mechanism to fully react and produce large \ce{NO} levels. Conversely, 1D flames simulate finite rate kinetics and realistic flow, within a limited space/time frame. The gas spends only milliseconds at peak temperatures, restricting the \ce{NO} formation and yielding drastically lower and realistic emissions.
\end{enumerate}